\Needspace{12\baselineskip}
\begin{center}\emph{Rigorous proof of the fundamental theorem.}\end{center}

\Article{135}
The preceding article did not prove the propositions of Article $133$, but showed that their truth depends on the truth of the fundamental theorem, which we assumed for the time being. The method of deduction itself, however, shows that they hold for the numbers $P,Q$ as long as the fundamental theorem holds when all the prime factors of these numbers are compared with one another, even if it does not hold generally. We now turn to the proof of the fundamental theorem itself, beginning with the following explanation.

\emph{We will say that the fundamental theorem is true up to a number $M$ if it holds for any two primes neither of which exceeds $M$.}

The same meaning applies when we say that the theorems of Articles $131,132,133$ are true up to some bound. It is easy to see that if the fundamental theorem is known to be true up to a bound, these propositions will also hold up to that same bound.

\SourceLeaf{110}\Article{136}
Checking the small numbers easily confirms the fundamental theorem and gives a bound up to which it certainly holds. We assume that this check has been made, but it does not matter how far we have taken it. It would even suffice to have confirmed it only up to the number $5$, and this requires just the single observation that $+5N3$, $\pm3N5$.

Now, if the fundamental theorem is not true generally, there will be some bound $T$ up to which it holds, but it will no longer hold up to the next number, $T+1$. This is the same as saying that there are two primes, the larger of which is $T+1$, whose comparison contradicts the fundamental theorem, while any other two primes, provided both are smaller than $T+1$, agree with it. So the propositions of Articles $131,132,133$ also hold up to $T$. We will now show that this supposition is impossible. We must distinguish the following cases according to the forms of $T+1$ and of the smaller prime, less than $T+1$, which we have supposed contradicts the theorem when compared with $T+1$. We denote that prime by $p$.

When both $T+1$ and $p$ have the form $4n+1$, the fundamental theorem could fail in two ways, namely if we had simultaneously \emph{either}

\[
\begin{array}{ll}&\pm pR(T+1)\quad\text{and}\quad\pm(T+1)Np\\\textit{or, simultaneously}&\pm pN(T+1)\quad\text{and}\quad\pm(T+1)Rp\end{array}
\]

When both $T+1$ and $p$ have the form $4n+3$, the fundamental theorem would fail if we had simultaneously \emph{either}

\[
\begin{array}{ll}&+pR(T+1)\quad\text{and}\quad-(T+1)Np\\\text{(or, equivalently}&-pN(T+1)\quad\text{and}\quad+(T+1)Rp)\\\textit{or}&+pN(T+1)\quad\text{and}\quad-(T+1)Rp\\\text{(or}&-pR(T+1)\quad\text{and}\quad+(T+1)Np)\end{array}
\]

\Needspace{8\baselineskip}
When $T+1$ has the form $4n+1$, but $p$ has the form $4n+3$, the fundamental theorem would fail if we had \emph{either}

\[
\begin{array}{ll}&\pm pR(T+1)\ \text{and}\ +(T+1)Np\ (\text{or}\ -(T+1)Rp)\\\textit{or}&\pm pN(T+1)\ \text{and}\ -(T+1)Np\ (\text{or}\ +(T+1)Rp)\end{array}
\]

\SourceLeaf{111}When $T+1$ has the form $4n+3$, but $p$ has the form $4n+1$, the fundamental theorem would fail if we had \emph{either}

\[
\begin{array}{ll}&+pR(T+1)\ (\text{or}\ -pN(T+1))\ \text{and}\ \pm(T+1)Np\\\textit{or}&+pN(T+1)\ (\text{or}\ -pR(T+1))\ \text{and}\ \pm(T+1)Rp\end{array}
\]

If we can prove that none of these eight cases can occur, we will also be certain that the truth of the fundamental theorem is not restricted by any bound. We now take up this task. But since some of these cases depend on others, we cannot keep the order in which we have listed them here.

\Needspace{9\baselineskip}
\Article{137}
\emph{First case. When $T+1$ has the form $4n+1$ $(=a)$, and $p$ has the same form, and moreover $\pm pRa$, we cannot have $\pm aNp$.} This was the first case above.

Let $+p\equiv e^2(\mathrm{mod.}\ a)$, with $e$ even and $<a$ (which can always be arranged). We must now distinguish two cases.

I.\quad When $e$ is not divisible by $p$, put $e^2=p+af$. Then $f$ is positive, of the form $4n+3$ (or of the form $B$), $<a$, and not divisible by $p$. Also $e^2\equiv p(\mathrm{mod.}\ f)$, \emph{that is,} $pRf$, and therefore, by Proposition $11$ of Article $132$, $\pm fRp$ (since $p,f<a$, these propositions hold for them). But we also have $afRp$, so $\pm aRp$ as well.

II.\quad When $e$ is divisible by $p$, put $e=gp$, and $e^2=p+aph$, or $pg^2=1+ah$. Then $h$ has the form $4n+3$ ($B$), and is coprime to $p$ and $g^2$. Also $pg^2Rh$, and therefore $pRh$, which by Proposition $11$ of Article $132$ gives $\pm hRp$. But we also have $-ahRp$, because $-ah\equiv1(\mathrm{mod.}\ p)$, so $\mp aRp$ as well.

\Needspace{9\baselineskip}
\Article{138}
\emph{Second case. When $T+1$ has the form $4n+1$ $(=a)$, $p$ the form $4n+3$, and $\pm pR(T+1)$, we cannot have $+(T+1)Np$, or $-(T+1)Rp$.} This was the fifth case above.

As above, let $e^2=p+fa$, with $e$ even and $<a$.

I.\quad When $e$ is not divisible by $p$, $f$ is not divisible by $p$ either. Moreover, $f$ is positive, of the form $4n+1$ (or $A$), and $<a$. We have $+pRf$, \SourceLeaf{112}and therefore $+fRp$ (Proposition 10 of Article 132). But we also have $+faRp$, so $+aRp$, or $-aNp$.

II. When $e$ is divisible by $p$, let $e=pg$, and $f=ph$. Then $g^2p=1+ha$. Now $h$ is positive, of the form $4n+3$ $(B)$, and coprime to $p$ and $g^2$. Also $+g^2pRh$, and therefore $+pRh$, which gives $-hRp$ (Proposition 13 of Article 132). But $-haRp$, so $+aRp$ and $-aNp$.

\Needspace{10\baselineskip}
\Article{139}
\textit{Third case. When $T+1$ has the form $4n+1$ $(=a)$, $p$ the same form, and $\pm pNa$, we cannot have $\pm aRp$.} (The second case above.)

Take a prime smaller than $a$ of which $+a$ is a nonresidue; we proved above that such primes exist (Articles 125, 129). We must consider two cases separately, according to whether this prime has the form $4n+1$ or $4n+3$, since it has not been proved that such primes exist in \textit{both} forms.

I. Let this prime have the form $4n+1$ and be $=a'$. Then $\pm a'Na$ (Article 131), and therefore $\pm a'pRa$. So let $e^2\equiv a'p\ (\mathrm{mod.}\ a)$, with $e$ even and $<a$. We must now distinguish four cases again.

1) When $e$ is divisible by neither $p$ nor $a'$, put $e^2=a'p\pm af$, taking the signs so that $f$ is positive. Then $f<a$, and it is coprime to $a'$ and $p$; for the upper sign it has the form $4n+3$, and for the lower sign the form $4n+1$. For brevity, let $[x,y]$ denote the number of prime factors of the number $y$ of which $x$ is a nonresidue. Then $a'pRf$, and therefore $[a'p,f]=0$. So $[f,a'p]$ is even (Propositions 1, 3 of Article 133), \textit{that is,} either $=0$ or $=2$. Thus $f$ is a residue of either both the numbers $a',p$ or neither. The former is impossible, since $\pm af$ is a residue of $a'$, and $\pm aNa'$ by hypothesis; this gives $\pm fNa'$. So $f$ must be a nonresidue of both the numbers $a',p$. But since $\pm afRp$, we have $\pm aNp$. \mbox{\textit{Q. E. D.}}

2) When $e$ is divisible by $p$ but not by $a'$, let $e=gp$, and $g^2p=a'\pm ah$, choosing the sign so that $h$ is positive. Then $h<a$, and it is coprime to $a',g$ and $p$; for the upper sign it has the form $4n+3$, and for the lower sign the form $4n+1$. From the equation $g^2p=a'\pm ah$, on multiplying by $p$ and $a'$, we easily deduce $pa'Rh\,.\,.\,.\,.\,.(\alpha)$; $\pm ahpRa'\,.\,.\,.\,.(\beta)$; $aa'hRp\,.\,.\,.\,.\,.(\gamma)$. From $(\alpha)$ follows $[pa',h]=0$, and therefore (Propositions 1, 3 of Article 133) $[h,pa']$ is even, \textit{that is,} \SourceLeaf{113}$h$ is a nonresidue of either both $p,a'$ or neither. \textit{In the former case,} $(\beta)$ gives $\pm apNa'$, and since $\pm aNa'$ by hypothesis, we have $\pm pRa'$. The fundamental theorem, which holds for the numbers $p,a'$ smaller than $T+1$, then gives $\pm a'Rp$. From this and $hNp$, by $(\gamma)$ we obtain $\pm aNp$. \mbox{\textit{Q. E. D.}} \textit{In the latter case,} $(\beta)$ gives $\pm apRa'$, then $\pm pNa'$, $\pm a'Np$, and finally, together with $hRp$, by $(\gamma)$ we obtain $\pm aNp$. \mbox{\textit{Q. E. D.}}

3) When $e$ is divisible by $a'$ but not by $p$, the proof proceeds almost exactly as in the preceding case. It will present no difficulty to anyone who has understood that case thoroughly.

4) When $e$ is divisible by both $a'$ and $p$, and therefore also by their product $a'p$ (we assume the numbers $a',p$ to be \textit{unequal}, since otherwise what we are trying to prove, namely $aNp$, would already be contained in the hypothesis $aNa'$), let $e=ga'p$, and $g^2a'p=1\pm ah$. Then $h<a$, and it is coprime to $a'$ and $p$; for the upper sign it has the form $4n+3$, and for the lower sign the form $4n+1$. It is easy to see that the following can be deduced from this equation: $a'pRh\,.\,.\,.\,.(\alpha)$; $\pm ahRa'\,.\,.\,.\,.(\beta)$; $\pm ahRp\,.\,.\,.\,.(\gamma)$. From $(\alpha)$, which agrees with $(\alpha)$ in (2), we find as there that either $hRp$, $hRa'$ hold simultaneously, or $hNp$, $hNa'$. But in the former case, $(\beta)$ would give $aRa'$, contrary to the hypothesis. So $hNp$, and therefore by $(\gamma)$ also $aNp$.

II. When this prime has the form $4n+3$, the proof is so similar to the preceding one that we thought it unnecessary to give it. For those eager to work it out for themselves (which we strongly recommend), we only note that, after reaching an equation such as $e^2=bp\pm af$ (where $b$ denotes that prime), it will make the argument clearer to consider the two signs separately.

\Needspace{8\baselineskip}
\Article{140}
\textit{Fourth case. When $T+1$ has the form $4n+1$ $(=a)$, $p$ the form $4n+3$, and $\pm pNa$, we cannot have $+aRp$, or $-aNp$.} (The sixth case above.)

Since its proof is exactly like the proof of the third case, we omit its proof too for brevity.

\SourceLeaf{114}\Needspace{9\baselineskip}
\Article{141}
\textit{Fifth case. When $T+1$ has the form $4n+3$ $(=b)$, $p$ the same form, and $+pRb$, or $-pNb$, we cannot have $+bRp$, or $-bNp$.} (The third case above.)

Let $p\equiv e^2\ (\mathrm{mod.}\ b)$, with $e$ even and $<b$.

I. When $e$ is not divisible by $p$, put $e^2=p+bf$. Then $f$ is positive, of the form $4n+3$, $<b$, and coprime to $p$. Also $pRf$, and therefore, by Proposition 13 of Article 132, $-fRp$. From this and $+bfRp$ we get $-bRp$, and therefore $+bNp$. \mbox{\textit{Q. E. D.}}

II. When $e$ is divisible by $p$, let $e=pg$, and $ggp=1+bh$. Then $h$ has the form $4n+1$ and is coprime to $p$; $p\equiv g^2p^2\ (\mathrm{mod.}\ h)$, and therefore $pRh$. This gives $+hRp$ (Proposition 10 of Article 132), which together with $-bhRp$ yields $-bRp$, or $+bNp$. \mbox{\textit{Q. E. D.}}

\Needspace{8\baselineskip}
\Article{142}
\textit{Sixth case. When $T+1$ has the form $4n+3$ $(=b)$, $p$ the form $4n+1$, and $pRb$, we cannot have $\pm bNp$.} (The seventh case above.)

We omit the proof, which is exactly like the preceding one.

\Needspace{9\baselineskip}
\Article{143}
\textit{Seventh case. When $T+1$ has the form $4n+3$ $(=b)$, $p$ the same form, and $+pNb$, or $-pRb$, we cannot have $+bNp$, or $-bRp$.} (The fourth case above.)

Let $-p\equiv e^2\ (\mathrm{mod.}\ b)$, with $e$ even and $<b$.

I. When $e$ is not divisible by $p$, let $-p=e^2-bf$. Then $f$ is positive, of the form $4n+1$, coprime to $p$, and smaller than $b$ (for $e$ is certainly no greater than $b-1$, and $p<b-1$, so $bf=e^2+p<b^2-b$, \textit{that is,} $f<b-1$). Also $-pRf$, which gives $+fRp$ (Proposition 10 of Article 132); from this and $+bfRp$ we get $+bRp$, or $-bNp$.

II. When $e$ is divisible by $p$, let $e=pg$, and $g^2p=-1+bh$. Then $h$ is positive, of the form $4n+3$, coprime to $p$, and $<b$. Also $-pRh$, which gives $+hRp$ (Proposition 14 of Article 132). From this and $bhRp$ follows $+bRp$, or $-bNp$. \mbox{\textit{Q. E. D.}}

\SourceLeaf{115}\Needspace{8\baselineskip}
\Article{144}
\textit{Eighth case. When $T+1$ has the form $4n+3$ $(=b)$, $p$ the form $4n+1$, and $+pNb$, or $-pRb$, we cannot have $\pm bRp$.} (The last case above.)

The proof proceeds just as in the preceding case.
